\documentclass[10pt,a4paper]{amsart}
\usepackage[margin=19mm]{geometry}
\usepackage{amsmath,amssymb,mathtools}
\usepackage[scr=rsfs,cal=euler]{mathalfa}
\usepackage{xfrac}
\usepackage[unicode,hidelinks,bookmarks=false]{hyperref}
\newcommand{\ie}{i.e.\ }
\newcommand{\cf}{cf.\ }
\newcommand{\resp}{respectively\ }
\newcommand{\Subsection}{\subsection}
\providecommand{\nobreakdash}{\nobreak\hbox{-}}
\setlength{\emergencystretch}{2em}
\multlinegap0pt
\usepackage{amssymb}
\usepackage{mathtools}
\usepackage{float}
\newtheorem{proposition}{Proposition}
\title{Path to homology of Yang-Baxter operators}
\author{Jozef H. Przytycki}
\date{Source: arXiv:2607.28626v1 (CC BY 4.0)}
\begin{document}
\maketitle

\noindent\emph{Source and licence.} Path to homology of Yang-Baxter operators. Version arXiv:2607.28626v1 (CC BY 4.0), distributed by arXiv under the Creative Commons Attribution 4.0 International licence, http://creativecommons.org/licenses/by/4.0/, as recorded in the arXiv licence metadata for this exact version and shown on the abstract page https://arxiv.org/abs/2607.28626v1. Full licence notice: \texttt{licenses/arxiv-2607.28626-CC-BY-4.0.txt}.\par\smallskip

\noindent\textit{Editorial scope.} Counted unit: the article's proposition Bin2 (source0.tex lines 1042--1045). The article's complete original proof of it is delivered with it (source0.tex lines 1046--1055, one \texttt{proof} environment of 10 lines). No mathematics is written, repaired or re-derived: the statement and the proof are the article's own lines, in the article's own notation, and no retained line is substituted or re-flowed.  No further source-local context is needed: every cross-reference of every retained line resolves inside the two delivered spans.  Nothing was omitted from any retained span, so the package carries no omission record.  No retained line contains a citation, so the wrapper carries no bibliography and no bibliography entry.  No retained line contains a figure environment, an image-inclusion command or drawing code, so no image is delivered and the package declares no asset dependency.  Editorial formatting only, in addition: an AMS \texttt{amsart} preamble that restates the article's own declarations and macros used by the retained lines, namely the environments \texttt{proposition} on separate counters, together with the standard packages those lines need.  The retained environments print in this document as Proposition 1 in order of appearance, whereas the article prints the same items with the numbering of its own full document.  The complete native source of the article as distributed in the arXiv e-print for this exact version is bundled unchanged as \texttt{sources/206430/source0.tex}.
\begin{proposition}\label{Bin2}
The set $Bin(X)$ forms a monoid with composition $*_1*_2$ of binary operations defined by $a(*_1*_2)b= (a*_1b)*_2b$ and the identity element 
$*_0$ defined by $a*_0b=a$.
\end{proposition}

\begin{proof} We check quickly that 
\begin{equation*}
a(**_0)b= (a*b)*_0b=a*b= (a*_0b)*b=a(*_0*)b,
\end{equation*}
thus $**_0=*=*_0*$. Associativity follows from associativity of composition of functions, we have:
\begin{equation*}
a(*_1*_2)*_3b=(a*_1*_2b)*_3b=((a*_1b)*_2b)*_3b=(a*_1b)*_2*_3b=a*_1(*_2*_3)b.  
\end{equation*}
Thus $(*_1*_2)*_3= *_1(*_2*_3)$, as needed.
\end{proof}

\end{document}
